% Generated draft: requires mathematical and source review.
\documentclass[11pt]{article}
\usepackage[margin=25mm]{geometry}
\usepackage{fontspec}
\setmainfont{Latin Modern Roman}
\usepackage{amsmath,amssymb,mathtools,unicode-math}
\setmathfont{Latin Modern Math}
\usepackage{enumitem,needspace,xurl,hyperref}
\setlength{\parindent}{0pt}
\setlength{\parskip}{5pt}
\setlength{\emergencystretch}{3em}
\newcommand{\N}{\mathbb{N}}
\newcommand{\Z}{\mathbb{Z}}
\newcommand{\Q}{\mathbb{Q}}
\newcommand{\R}{\mathbb{R}}
\newcommand{\C}{\mathbb{C}}
\newcommand{\F}{\mathbb{F}}
\newcommand{\A}{\mathbb{A}}
\newcommand{\PP}{\mathbb{P}}
\newcommand{\E}{\mathbb{E}}
\newcommand{\eps}{\varepsilon}
\newcommand{\dd}{\,\mathrm{d}}
\begin{document}
\section*{UnsolvedMath ID: 30000400}
\textbf{Source problem code:} catalogue entry
\subsection{Definitions and mathematical statement}
The exact source question is: for a specified random microstructure and prescribed accuracy, at what mesoscale is a representative volume element attained?

A representative-volume criterion compares an apparent quantity $A_L(\omega)$ computed on a sample of linear scale $L$ with a prescribed effective or ensemble target $A_{\mathrm{eff}}$. Once the material law, sample boundary conditions, observable and norm are fixed, one possible probabilistic accuracy definition is
\[
 \mathbb P\bigl(\|A_L-A_{\mathrm{eff}}\|>\varepsilon\bigr)\le\delta,
 \qquad\varepsilon>0,\quad0<\delta<1.
\]
A persistent threshold would seek $L_0$ for which this holds for every admissible $L\ge L_0$. Mean-square error, bias bounds or confidence intervals are different possible accuracy conventions; the source does not choose one. The displayed inequality therefore illustrates what must be quantified, not an added source requirement or a claimed universal law.

The extract supplies no probability law, correlation/mixing assumptions, microstructural length, effective observable, boundary-condition prescription or error/confidence convention. It also does not guarantee an effective limit under unprovided stationarity/ergodicity hypotheses. Consequently a universal numerical mesoscale or convergence rate cannot be inferred. The task is to determine the needed representative scale for the actually specified model and tolerance, not to assume that one scale works for every random material.

The literature assessment checked on 21 August 2026 explicitly calls the size-threshold question under-specified and says its resolution cannot be verified from the extraction. The registry's open label is retained as dated curation, not an independently verified unresolved theorem.
\subsection{Short English statement}
For a given random material and accuracy criterion, determine how large a sample must be to represent the effective behavior. The sample model, observable, boundary conditions and meaning of accuracy are missing from this extraction.
\subsection{Sources}
\begin{itemize}
\item \emph{Mechanics of Materials} (2006), Oberwolfach Reports, DOI 10.4171/owr/2006/04, pp. 187--252, \url{https://doi.org/10.4171/owr/2006/04}.
\item The sole literature source in the assessment checked 21 August 2026 is the original report, \url{https://doi.org/10.4171/owr/2006/04}; the assessment explicitly reports missing defining context.
\item Local UnsolvedMath registry, permanent ID $30000400$, exact source code \texttt{catalogue entry}, created 25 August 2026 and updated 16 September 2026. The statement status is exact and the broad sentence was retained byte-for-byte.
\end{itemize}
\end{document}
